\begin{frame}{두 방법은 왜 같은 지점에 도달할까}
  \begin{block}{핵심}
    두 방법이 결국 같은 지점에 도달하는 것은 \textbf{우연이 아니다} ---
    $J(\boldsymbol{w})$ 가 \kb{볼록함수}라서 유일한 전역 최솟값을 가지므로,
    ``미분해서 한 번에 푼다''와 ``조금씩 내려간다''는 서로 다른
    두 길이 \textbf{같은 목적지}로 이어지는 것뿐이다.
  \end{block}
\end{frame}
